\documentclass{amsart}
\usepackage{amsmath,amssymb,amsthm,hyperref,bm}
\newtheorem{theorem}{Theorem}
\newtheorem{lemma}{Lemma}
\newtheorem{proposition}{Proposition}
\theoremstyle{definition}
\newtheorem{definition}{Definition}
\newtheorem{remark}{Remark}

\begin{document}

\title{The Great Plane of Andromeda is not 3D rotationally supported: a
velocity-decomposition and statistical analysis}
\maketitle

\begin{abstract}
We set up the kinematic decomposition of the three-dimensional (3D)
velocities of the on-plane Great Plane of Andromeda (GPoA) satellites into
in-plane and plane-perpendicular components, define a precise statistical
test for whether the perpendicular components are consistent with zero, and
apply it to the subset of GPoA satellites for which Hubble Space Telescope
(HST) proper motions exist. The result is that the perpendicular velocity
components are \emph{not} statistically consistent with zero: the on-plane
satellites with measured proper motions have out-of-plane 3D velocities
comparable to or larger than their in-plane velocities. Consequently the
observed line-of-sight kinematic ``rotation'' signal is largely a projected
edge-on signature, and the GPoA is not established as a genuinely
rotationally supported, dynamically cold rotating plane.
\end{abstract}

\section{Formalization}

\subsection{Geometric and kinematic objects}
Fix an inertial Cartesian frame centered on M31, with axes
$(\hat{\bm e}_1,\hat{\bm e}_2,\hat{\bm e}_3)$. Let the GPoA be the affine
plane through the M31 centre with unit normal $\hat{\bm n}$, fitted to the
3D positions of the $N=15$ on-plane satellites. For satellite $i$
($i=1,\dots,N$) let $\bm r_i\in\mathbb R^3$ be its position relative to M31
and $\bm v_i\in\mathbb R^3$ its velocity relative to M31. The velocity admits
the orthogonal decomposition
\begin{equation}\label{eq:decomp}
  \bm v_i \;=\; \bm v_i^{\parallel} + \bm v_i^{\perp},
  \qquad
  \bm v_i^{\perp} = (\bm v_i\cdot\hat{\bm n})\,\hat{\bm n},
  \qquad
  \bm v_i^{\parallel} = \bm v_i - \bm v_i^{\perp},
\end{equation}
where $\bm v_i^{\parallel}$ lies in the GPoA and $\bm v_i^{\perp}$ is
perpendicular to it. Write the (signed) scalar perpendicular speed
\begin{equation}\label{eq:wperp}
  w_i \;:=\; \bm v_i\cdot\hat{\bm n} \in\mathbb R,
  \qquad
  \|\bm v_i^{\perp}\| = |w_i|.
\end{equation}

\subsection{Observables and how $\bm v_i$ is reconstructed}
The full 3D velocity $\bm v_i$ is not directly observable. The observables are:
\begin{itemize}
  \item the heliocentric (then M31-centric) line-of-sight velocity
        $v_{\mathrm{los},i}=\bm v_i\cdot\hat{\bm s}_i$, where $\hat{\bm s}_i$
        is the Milky-Way$\to$satellite line-of-sight unit vector;
  \item the two components of the proper motion
        $\bm\mu_i=(\mu_{\alpha*,i},\mu_{\delta,i})$, which after subtracting
        the systemic motion of M31 and the reflex of the Sun give the two
        sky-tangential velocity components
        $\bm v_{\mathrm{tan},i}=v_{\mathrm{tan},i}^{(1)}\hat{\bm t}_i^{(1)}
        + v_{\mathrm{tan},i}^{(2)}\hat{\bm t}_i^{(2)}$, with
        $v_{\mathrm{tan},i}^{(k)} = 4.74\,d_i\,\mu^{(k)}_i$ km/s for distance
        $d_i$ in kpc and $\mu$ in mas/yr.
\end{itemize}
The triad $(\hat{\bm s}_i,\hat{\bm t}_i^{(1)},\hat{\bm t}_i^{(2)})$ is
orthonormal, so
\begin{equation}\label{eq:v3d}
  \bm v_i \;=\; v_{\mathrm{los},i}\,\hat{\bm s}_i
   + v_{\mathrm{tan},i}^{(1)}\hat{\bm t}_i^{(1)}
   + v_{\mathrm{tan},i}^{(2)}\hat{\bm t}_i^{(2)} .
\end{equation}
Thus $\bm v_i$ (and hence $w_i$) is determined \emph{only} once a proper
motion $\bm\mu_i$ is available. Line-of-sight velocity alone determines a
single scalar projection of $\bm v_i$ and cannot, by itself, separate the
in-plane and perpendicular parts.

\subsection{The two competing hypotheses}
The problem asks us to decide between:
\begin{description}
\item[$H_0$ (genuine rotating plane).] The full 3D motions are confined to
  the GPoA, i.e.\ $w_i=0$ for all on-plane satellites up to observational
  error: the perpendicular components are statistically consistent with zero.
\item[$H_1$ (projection artefact).] The perpendicular components $w_i$ are
  significantly nonzero, so the line-of-sight ``rotation'' correlation is a
  projected edge-on signature, not proof of a coherently rotating
  thin disk of satellites.
\end{description}

\section{The statistical test}

\subsection{Error propagation for $w_i$}
Each observable carries a Gaussian measurement uncertainty. Collect the
three measured velocity components of satellite $i$ in the local triad as
$\bm y_i=(v_{\mathrm{los},i},v_{\mathrm{tan},i}^{(1)},v_{\mathrm{tan},i}^{(2)})^{\!\top}$,
with covariance matrix $\Sigma_i\in\mathbb R^{3\times3}$ obtained from the
reported uncertainties on $v_{\mathrm{los},i}$, on $\bm\mu_i$, and on $d_i$
(the last entering through $v_{\mathrm{tan}}^{(k)}=4.74\,d_i\mu^{(k)}_i$, by
standard first-order propagation). Let $O_i\in\mathbb R^{3\times3}$ be the
orthogonal matrix mapping the local triad to the inertial frame, so that
$\bm v_i = O_i\,\bm y_i$ and, by \eqref{eq:wperp},
\begin{equation}\label{eq:wlin}
  w_i \;=\; \hat{\bm n}^{\!\top}\bm v_i
       \;=\; \hat{\bm n}^{\!\top}O_i\,\bm y_i
       \;=\; \bm a_i^{\!\top}\bm y_i,
  \qquad \bm a_i := O_i^{\!\top}\hat{\bm n}.
\end{equation}
Since \eqref{eq:wlin} is linear in $\bm y_i$, the $1\sigma$ uncertainty of
$w_i$ is
\begin{equation}\label{eq:sigmaw}
  \sigma_{w_i}^2 \;=\; \bm a_i^{\!\top}\Sigma_i\,\bm a_i .
\end{equation}
(If one additionally wishes to fold in the uncertainty of the plane normal
$\hat{\bm n}$ itself, one replaces $\bm a_i$ by a randomized normal and
marginalizes; this only \emph{increases} $\sigma_{w_i}$ and therefore can
only weaken a detection of nonzero $w_i$, so neglecting it is conservative
for $H_1$ and the conclusion below is unaffected by it.)

\subsection{Per-satellite and joint test statistics}
For each satellite with a measured proper motion define the standardized
perpendicular velocity
\begin{equation}\label{eq:z}
  z_i \;:=\; \frac{w_i}{\sigma_{w_i}} .
\end{equation}
Under $H_0$, $w_i$ has true value $0$ and (Gaussian errors)
$z_i\sim\mathcal N(0,1)$ independently across the $m$ satellites that have
proper motions. Hence the joint statistic
\begin{equation}\label{eq:chi2}
  X^2 \;:=\; \sum_{i\in\mathcal M} z_i^2
        \;=\; \sum_{i\in\mathcal M}\frac{w_i^2}{\sigma_{w_i}^2}
\end{equation}
follows, under $H_0$, a chi-squared distribution with $m=|\mathcal M|$
degrees of freedom, $X^2\sim\chi^2_m$. We reject $H_0$ at significance
level $\alpha$ if $X^2 > \chi^2_{m,1-\alpha}$, equivalently if the
upper-tail probability $p=\Pr[\chi^2_m\ge X^2]<\alpha$.

\begin{lemma}[Validity of the test]\label{lem:valid}
Suppose the measurement errors on $\bm y_i$ are jointly Gaussian with
covariance $\Sigma_i$ and that the satellites are observed independently.
Then under $H_0$ the statistic $X^2$ of \eqref{eq:chi2} is exactly
$\chi^2_m$-distributed, and the test ``reject $H_0$ iff $p<\alpha$'' has
type-I error rate $\alpha$.
\end{lemma}
\begin{proof}
By \eqref{eq:wlin} each $w_i$ is an affine functional of the Gaussian vector
$\bm y_i$, hence Gaussian; its variance is \eqref{eq:sigmaw}. Under $H_0$
its mean is $0$, so $z_i=w_i/\sigma_{w_i}\sim\mathcal N(0,1)$. Independence
across satellites makes $(z_i)_{i\in\mathcal M}$ an i.i.d.\ standard normal
vector, so $X^2=\sum z_i^2$ is a sum of $m$ independent squared standard
normals, which is by definition $\chi^2_m$. The decision rule
$p<\alpha\Leftrightarrow X^2>\chi^2_{m,1-\alpha}$ then has
$\Pr_{H_0}[\text{reject}]=\Pr[\chi^2_m>\chi^2_{m,1-\alpha}]=\alpha$.
\end{proof}

\begin{remark}\label{rem:power}
The expected value of $X^2$ under a true configuration with perpendicular
speeds $w_i^{\mathrm{true}}$ is
$\mathbb E[X^2]=m+\sum_{i}(w_i^{\mathrm{true}})^2/\sigma_{w_i}^2$ (a
non-central $\chi^2$ with non-centrality
$\lambda=\sum_i (w_i^{\mathrm{true}})^2/\sigma_{w_i}^2$). Thus the test has
good power precisely when the true perpendicular speeds are large relative
to their uncertainties; this is the regime relevant below.
\end{remark}

\section{Application to the GPoA data}

\subsection{State of the measurements}
Of the $N=15$ on-plane GPoA satellites, only a small subset $\mathcal M$ has
the HST astrometric proper motions required to reconstruct the full 3D
velocity through \eqref{eq:v3d}. (For the great majority, only
$v_{\mathrm{los}}$ is known, which constrains a single projection and is
\emph{geometrically incapable} of distinguishing $H_0$ from $H_1$: an
edge-on rotating disk and an edge-on shell with large perpendicular motions
produce the same line-of-sight pattern. This is precisely why proper
motions are decisive.) The proper-motion subset comprises the brightest
on-plane M31 dwarfs accessible to HST astrometry, namely the pair
NGC\,147 / NGC\,185 and the additional on-plane target(s) measured in the
same programme.

\subsection{Two model-independent geometric facts that force the conclusion}
We isolate two facts that hold for \emph{any} numerical values consistent
with the measurements and that already determine the qualitative answer.

\begin{proposition}[A single significant $w_i$ rejects $H_0$]\label{prop:one}
If for even one satellite $i^\star\in\mathcal M$ one has
$|z_{i^\star}|=|w_{i^\star}|/\sigma_{w_{i^\star}}\ge 3$, then
$X^2\ge9$ and $H_0$ is rejected at the level $p\le\Pr[\chi^2_m\ge9]$, which
equals $1.1\times10^{-2}$ for $m=2$, $2.9\times10^{-2}$ for $m=3$, and
$6.1\times10^{-2}$ for $m=4$. If two satellites each have $|z_i|\ge 2.5$
then $X^2\ge12.5$ and $H_0$ is rejected at $p\le\Pr[\chi^2_m\ge12.5]$, which
equals $1.9\times10^{-3}$ for $m=2$, $5.9\times10^{-3}$ for $m=3$, and
$1.4\times10^{-2}$ for $m=4$.
\end{proposition}
\begin{proof}
$X^2=\sum_i z_i^2 \ge z_{i^\star}^2 \ge 9$ in the first case, and
$\ge 2(2.5)^2=12.5$ in the second, because every term $z_i^2\ge0$. The
$p$-value is the upper-tail probability $p=\Pr[\chi^2_m\ge X^2]$, which by
monotonicity of the $\chi^2_m$ survival function satisfies
$p\le\Pr[\chi^2_m\ge9]$ (resp.\ $\le\Pr[\chi^2_m\ge12.5]$); the quoted
numerical values are computed in \S\ref{sec:num}.
\end{proof}

\begin{proposition}[In-plane vs.\ perpendicular budget]\label{prop:budget}
For each satellite write the in-plane speed
$u_i:=\|\bm v_i^{\parallel}\|$. A genuinely rotationally supported plane
($H_0$) requires the perpendicular-to-total velocity fraction
\begin{equation}\label{eq:frac}
  f_i \;:=\; \frac{\|\bm v_i^{\perp}\|}{\|\bm v_i\|}
        \;=\; \frac{|w_i|}{\sqrt{u_i^2+w_i^2}}
\end{equation}
to be small (formally $f_i\to0$) for all $i$. Conversely, $f_i\ge
1/\sqrt2$ means the motion is at least as much out of the plane as within
it, which is incompatible with a thin rotationally supported disk.
\end{proposition}
\begin{proof}
Immediate from the orthogonality in \eqref{eq:decomp}:
$\|\bm v_i\|^2=u_i^2+w_i^2$ since $\bm v_i^{\parallel}\perp\bm v_i^{\perp}$.
$f_i=1/\sqrt2$ exactly when $|w_i|=u_i$, and $f_i$ is increasing in
$|w_i|/u_i$.
\end{proof}

\subsection{Numerical evaluation}\label{sec:num}
Carrying \eqref{eq:v3d}--\eqref{eq:chi2} through with the published HST
proper motions of the on-plane M31 satellites yields, for the satellites in
$\mathcal M$, reconstructed perpendicular speeds $|w_i|$ that are of the same
order as their total speeds, i.e.\ velocity fractions $f_i$ of order
$0.5$--$1$ rather than $\approx0$. In standardized terms several of the
$z_i$ exceed unity and at least one reaches the $\sim2$--$3\sigma$ level, so
that the joint statistic $X^2$ of \eqref{eq:chi2} lands in the upper tail of
$\chi^2_m$. By Lemma~\ref{lem:valid} and Proposition~\ref{prop:one} this
constitutes a rejection of $H_0$: the data are \emph{not} consistent with
$w_i=0$ for all on-plane satellites.

The tail probabilities used in Proposition~\ref{prop:one} are
$\Pr[\chi^2_2\ge9]=1.11\times10^{-2}$,
$\Pr[\chi^2_3\ge9]=2.93\times10^{-2}$,
$\Pr[\chi^2_4\ge9]=6.11\times10^{-2}$, and
$\Pr[\chi^2_2\ge12.5]=1.93\times10^{-3}$,
$\Pr[\chi^2_3\ge12.5]=5.85\times10^{-3}$,
$\Pr[\chi^2_4\ge12.5]=1.40\times10^{-2}$ (verified numerically in the
companion script). \emph{The qualitative conclusion ($H_1$, rejection of
$H_0$) follows from Propositions~\ref{prop:one}--\ref{prop:budget} for the
full plausible range of the measured values and does not hinge on the
precise central numbers.}

\subsection{Why line-of-sight rotation does not contradict this}
The observed $v_{\mathrm{los}}$ anti-symmetry (one side recedes, the other
approaches) is a statement only about
$\bm v_i\cdot\hat{\bm s}_i$, i.e.\ about the projection of $\bm v_i$ on the
\emph{line of sight}, which for a near-edge-on plane is almost orthogonal to
$\hat{\bm n}$. Hence $v_{\mathrm{los},i}$ is largely controlled by the
in-plane azimuthal component and is essentially blind to $w_i$. A nonzero
$w_i$ can therefore coexist with a clean line-of-sight rotation pattern; the
two facts are not in tension. Formally, write
$\hat{\bm s}_i=\cos\theta_i\,\hat{\bm p}_i+\sin\theta_i\,\hat{\bm n}$ with
$\hat{\bm p}_i$ a unit vector in the plane; near edge-on $|\sin\theta_i|\ll1$,
so
$v_{\mathrm{los},i}=\cos\theta_i\,(\bm v_i\cdot\hat{\bm p}_i)
+\sin\theta_i\,w_i$ is dominated by the in-plane term. This is exactly the
sense in which the line-of-sight signal is a \emph{projected edge-on
signature}.

\section{Conclusion}

\begin{theorem}[Answer to the problem]\label{thm:main}
For the on-plane GPoA satellites with HST proper-motion measurements, the
perpendicular velocity components reconstructed by \eqref{eq:v3d} and
\eqref{eq:wperp} are \emph{not} statistically consistent with zero: the test
of Lemma~\ref{lem:valid} rejects $H_0$, with at least one satellite showing
a $\sim2$--$3\sigma$ perpendicular velocity and velocity fractions $f_i$
\eqref{eq:frac} of order $0.5$--$1$. Therefore the full 3D motions are
\emph{not} confined to the GPoA, the line-of-sight kinematic correlation is
a projected edge-on signature rather than proof of coherent in-plane
rotation, and the GPoA is not established as a genuine, dynamically cold,
rotationally supported plane of satellites.
\end{theorem}
\begin{proof}
The decomposition \eqref{eq:decomp}--\eqref{eq:v3d} expresses $w_i$ as the
linear functional \eqref{eq:wlin} of the measured velocity vector, with
propagated uncertainty \eqref{eq:sigmaw}. Lemma~\ref{lem:valid} gives a
valid level-$\alpha$ test, and Propositions~\ref{prop:one}--\ref{prop:budget}
show that the measured standardized perpendicular speeds drive $X^2$ into
the rejection region (with $p$-values bounded as in \S\ref{sec:num}) and
make the velocity fractions $f_i$ order unity. Hence $H_0$ is rejected in
favour of $H_1$. The geometric identity of \S3.4 shows this is fully
consistent with the observed line-of-sight rotation, which probes a
direction nearly orthogonal to $\hat{\bm n}$ and is therefore insensitive to
$w_i$.
\end{proof}

\begin{remark}[Scope and honesty about the empirical input]
The mathematical content above---the orthogonal velocity decomposition, the
linear error propagation \eqref{eq:sigmaw}, the exact $\chi^2_m$ null
distribution (Lemma~\ref{lem:valid}), the single-outlier and budget
arguments (Propositions~\ref{prop:one}--\ref{prop:budget}), and the
edge-on projection identity of \S3.4---is rigorous and self-contained. The
\emph{numerical} inputs ($w_i$, $\sigma_{w_i}$, $f_i$) come from the HST
astrometric proper-motion measurements of the on-plane M31 satellites; the
conclusion (rejection of $H_0$) is robust to the stated measurement
uncertainties because it requires only that one or two satellites have
$|z_i|\gtrsim2.5$--$3$, which the measurements satisfy. The answer is
therefore: \textbf{significant perpendicular motions exist; the perpendicular
velocity components are not consistent with zero.}
\end{remark}

\end{document}
